% META: {"id": "TEXP-ARI-EX-023", "chapitre": "TEXP-ARITHMETIQUE", "type_objet": "exercice", "capacites": ["TEXP-ARI-C7"], "capacites_codes": ["C7"], "methodes": ["M7"], "parcours": 2, "competences": ["raisonner"], "duree_min": 15, "mode_creation": "ex_nihilo", "sources_inspiration": [], "fichier_tex": "chapitres/TEXP-ARITHMETIQUE/exercices/TEXP-ARI-EX-023.tex", "corrige_tex": "chapitres/TEXP-ARITHMETIQUE/corriges/TEXP-ARI-CO-023.tex", "status": "generated"}
% BEGIN-VERIFY
% import math
% def premier(n):
%     return n > 1 and all(n % k for k in range(2, int(math.isqrt(n)) + 1))
% premiers = [p for p in range(2, 60) if premier(p)]
% for r in range(1, 7):
%     produit = 1
%     for p in premiers[:r]:
%         produit *= p
%     N = produit + 1
%     for p in premiers[:r]:
%         assert N % p != 0
% assert 2 * 3 * 5 * 7 + 1 == 211 and premier(211)
% assert 2 * 3 * 5 * 7 * 11 * 13 + 1 == 30031 == 59 * 509
% assert not premier(30031)
% END-VERIFY
\begin{exercice}{TEXP-ARI-EX-023}{2}{15}
On reprend la démonstration d'Euclide de l'infinitude des nombres premiers.
\begin{enumerate}
  \item Supposons qu'il n'existe qu'un nombre fini de premiers
        $p_1,\dots,p_r$. On pose $N=p_1p_2\cdots p_r+1$. Montrer qu'aucun $p_i$
        ne divise $N$.
  \item En déduire une contradiction, et conclure.
  \item Appliquer la construction aux quatre premiers nombres premiers : que
        vaut $N$ ? Est-il premier ?
\end{enumerate}
\end{exercice}
